Un des objectifs principaux en physique de la matière condensée vise à améliorer la compréhension que nous avons des matériaux, mais surtout à se servir de l’étude de leurs propriétés pour en découvrir de nouveaux. Avec le temps, ces derniers deviennent de plus en plus exotiques, notamment pour se plier aux exigences technologiques. Cela encourage naturellement les physiciens et physiciennes à redoubler d'efforts pour trouver des théories capables d'expliquer ces matériaux.

Ce mémoire se divise en trois chapitres et deux annexes. Le chapitre~\ref{chap: outils_theoriques} établit les bases théoriques nécessaires à la compréhension de la CDMFT. On y présente notamment la base de Fock, le modèle de Hubbard, le formalisme des fonctions de Green et une introduction à la théorie des groupes. Le chapitre~\ref{chap: methodes_numeriques} détaille les méthodes numériques en expliquant comment, à partir de la CPT, il est possible de définir la CDMFT. Le concept des sous-bains est ensuite intégré à la méthode et on expose comment la théorie des groupes peut s’avérer un outil puissant dans le contexte de la méthode SB-CDMFT. Le chapitre~\ref{chap: resultats} présente des résultats d'étalonnage, validant l'approche sur des systèmes de référence, et introduit une publication scientifique dans laquelle sont discutés les résultats les plus importants. On y trouve des systèmes résolus avec la SB-CDMFT dont la référence CDMFT n'existe pas par manque de ressources computationnelles, une des victoires de la méthode.

L'annexe~\ref{chap: ph-symmetry} contient une dérivation analytique de la symétrie particule-trou pour le modèle d'impureté d'Anderson et son effet sur la fonction de Green, permettant d'expliquer les symétries observées dans les résultats. Enfin, l'annexe~\ref{chap: article} est une copie d'un article publié dans la revue \textit{Physical Review B} présentant en détails la méthode SB-CDMFT.
